% Seminar 3 - Page 2. Fragment for
% % Seminar 3 - Page 2. Fragment for
% % Seminar 3 - Page 2. Fragment for
% % Seminar 3 - Page 2. Fragment for
% \input{seminars/pages/seminar3-pagina2-en}; no preamble or forced page break.
% EDITORIAL NOTES:
% - This page continues Problem 2(b) from Seminar 3 - Page 1.
% - The handwritten notation varies between P_y/P_z and p_y/p_z;
%   the projection maps are typeset consistently as p_Y and p_Z.
% - The diagram below rearranges the handwritten projection sketch for
%   clarity; the two component maps and the projections are retained.
% - A short aside seems to read: fix a in Z and define i_Y(y)=(y,a).
%   This is included, although it is not used in the following proof.
% - Crossed-out calculations near the two composition identities are
%   omitted; all legible final calculations are retained.
% - The manuscript occasionally writes an identity without a domain;
%   the domains are made explicit in the displayed proof.

\subsection*{Seminar 3 - Page 2}

\noindent\textbf{(b) Continued.} We have the correspondence
\[
  \operatorname{Hom}(X,Y\times Z)
  \;\overset{\alpha}{\underset{\beta}{\rightleftarrows}}\;
  \operatorname{Hom}(X,Y)\times\operatorname{Hom}(X,Z).
\]
Let $p_Y:Y\times Z\to Y$ and $p_Z:Y\times Z\to Z$ be the
canonical projections:
\[
 p_Y(y,z)=y,\qquad p_Z(y,z)=z.
\]
The diagram representing the projections and the component maps is
\[
\begin{tikzcd}[column sep=large, row sep=large]
 & X \arrow[dl, "h"'] \arrow[dr, "g"]
     \arrow[d, dashed, "{\beta(g,h)}"] & \\
 Z & Y\times Z \arrow[l, "p_Z"] \arrow[r, "p_Y"'] & Y
\end{tikzcd}
\]
For $f:X\to Y\times Z$ we define
\[
  \alpha(f)=(p_Y\circ f,\,p_Z\circ f).
\]
Conversely, for $g:X\to Y$ and $h:X\to Z$, set
\[
  \beta(g,h)(x)=(g(x),h(x)),\qquad x\in X.
\]

\medskip
\noindent\emph{Aside in the manuscript.} Fix $a\in Z$ and consider
\[
 \iota_Y:Y\longrightarrow Y\times Z,\qquad
 \iota_Y(y)=(y,a).
\]

\medskip
\noindent We prove first that $\alpha\circ\beta=\operatorname{id}$.
For every $g:X\to Y$ and $h:X\to Z$,
\[
  \alpha\bigl(\beta(g,h)\bigr)=(g,h).
\]
Indeed,
\[
\begin{aligned}
 \bigl(p_Y\circ\beta(g,h)\bigr)(x)
   &=p_Y\bigl(\beta(g,h)(x)\bigr) \\
   &=p_Y(g(x),h(x))=g(x),
 \qquad \forall x\in X,\\[3pt]
 \bigl(p_Z\circ\beta(g,h)\bigr)(x)
   &=p_Z\bigl(\beta(g,h)(x)\bigr) \\
   &=p_Z(g(x),h(x))=h(x),
 \qquad \forall x\in X.
\end{aligned}
\]
Consequently,
\[
  (\alpha\circ\beta)(g,h)
    =(p_Y\circ\beta(g,h),\,p_Z\circ\beta(g,h))
    =(g,h).
\]

\medskip
\noindent Conversely, we show that $\beta\circ\alpha=\operatorname{id}$.
Let $f:X\to Y\times Z$. Then
\[
 (\beta\circ\alpha)(f)=f.
\]
For $x\in X$, write $f(x)=(y,z)$. We have
\[
\begin{aligned}
 \beta\bigl(\alpha(f)\bigr)(x)
   &=\beta(p_Y\circ f,\,p_Z\circ f)(x)\\
   &=\bigl((p_Y\circ f)(x),\,(p_Z\circ f)(x)\bigr)\\
   &=\bigl(p_Y(y,z),\,p_Z(y,z)\bigr)\\
   &=(y,z)=f(x).
\end{aligned}
\]
Thus the two correspondences are mutually inverse.
; no preamble or forced page break.
% EDITORIAL NOTES:
% - This page continues Problem 2(b) from Seminar 3 - Page 1.
% - The handwritten notation varies between P_y/P_z and p_y/p_z;
%   the projection maps are typeset consistently as p_Y and p_Z.
% - The diagram below rearranges the handwritten projection sketch for
%   clarity; the two component maps and the projections are retained.
% - A short aside seems to read: fix a in Z and define i_Y(y)=(y,a).
%   This is included, although it is not used in the following proof.
% - Crossed-out calculations near the two composition identities are
%   omitted; all legible final calculations are retained.
% - The manuscript occasionally writes an identity without a domain;
%   the domains are made explicit in the displayed proof.

\subsection*{Seminar 3 - Page 2}

\noindent\textbf{(b) Continued.} We have the correspondence
\[
  \operatorname{Hom}(X,Y\times Z)
  \;\overset{\alpha}{\underset{\beta}{\rightleftarrows}}\;
  \operatorname{Hom}(X,Y)\times\operatorname{Hom}(X,Z).
\]
Let $p_Y:Y\times Z\to Y$ and $p_Z:Y\times Z\to Z$ be the
canonical projections:
\[
 p_Y(y,z)=y,\qquad p_Z(y,z)=z.
\]
The diagram representing the projections and the component maps is
\[
\begin{tikzcd}[column sep=large, row sep=large]
 & X \arrow[dl, "h"'] \arrow[dr, "g"]
     \arrow[d, dashed, "{\beta(g,h)}"] & \\
 Z & Y\times Z \arrow[l, "p_Z"] \arrow[r, "p_Y"'] & Y
\end{tikzcd}
\]
For $f:X\to Y\times Z$ we define
\[
  \alpha(f)=(p_Y\circ f,\,p_Z\circ f).
\]
Conversely, for $g:X\to Y$ and $h:X\to Z$, set
\[
  \beta(g,h)(x)=(g(x),h(x)),\qquad x\in X.
\]

\medskip
\noindent\emph{Aside in the manuscript.} Fix $a\in Z$ and consider
\[
 \iota_Y:Y\longrightarrow Y\times Z,\qquad
 \iota_Y(y)=(y,a).
\]

\medskip
\noindent We prove first that $\alpha\circ\beta=\operatorname{id}$.
For every $g:X\to Y$ and $h:X\to Z$,
\[
  \alpha\bigl(\beta(g,h)\bigr)=(g,h).
\]
Indeed,
\[
\begin{aligned}
 \bigl(p_Y\circ\beta(g,h)\bigr)(x)
   &=p_Y\bigl(\beta(g,h)(x)\bigr) \\
   &=p_Y(g(x),h(x))=g(x),
 \qquad \forall x\in X,\\[3pt]
 \bigl(p_Z\circ\beta(g,h)\bigr)(x)
   &=p_Z\bigl(\beta(g,h)(x)\bigr) \\
   &=p_Z(g(x),h(x))=h(x),
 \qquad \forall x\in X.
\end{aligned}
\]
Consequently,
\[
  (\alpha\circ\beta)(g,h)
    =(p_Y\circ\beta(g,h),\,p_Z\circ\beta(g,h))
    =(g,h).
\]

\medskip
\noindent Conversely, we show that $\beta\circ\alpha=\operatorname{id}$.
Let $f:X\to Y\times Z$. Then
\[
 (\beta\circ\alpha)(f)=f.
\]
For $x\in X$, write $f(x)=(y,z)$. We have
\[
\begin{aligned}
 \beta\bigl(\alpha(f)\bigr)(x)
   &=\beta(p_Y\circ f,\,p_Z\circ f)(x)\\
   &=\bigl((p_Y\circ f)(x),\,(p_Z\circ f)(x)\bigr)\\
   &=\bigl(p_Y(y,z),\,p_Z(y,z)\bigr)\\
   &=(y,z)=f(x).
\end{aligned}
\]
Thus the two correspondences are mutually inverse.
; no preamble or forced page break.
% EDITORIAL NOTES:
% - This page continues Problem 2(b) from Seminar 3 - Page 1.
% - The handwritten notation varies between P_y/P_z and p_y/p_z;
%   the projection maps are typeset consistently as p_Y and p_Z.
% - The diagram below rearranges the handwritten projection sketch for
%   clarity; the two component maps and the projections are retained.
% - A short aside seems to read: fix a in Z and define i_Y(y)=(y,a).
%   This is included, although it is not used in the following proof.
% - Crossed-out calculations near the two composition identities are
%   omitted; all legible final calculations are retained.
% - The manuscript occasionally writes an identity without a domain;
%   the domains are made explicit in the displayed proof.

\subsection*{Seminar 3 - Page 2}

\noindent\textbf{(b) Continued.} We have the correspondence
\[
  \operatorname{Hom}(X,Y\times Z)
  \;\overset{\alpha}{\underset{\beta}{\rightleftarrows}}\;
  \operatorname{Hom}(X,Y)\times\operatorname{Hom}(X,Z).
\]
Let $p_Y:Y\times Z\to Y$ and $p_Z:Y\times Z\to Z$ be the
canonical projections:
\[
 p_Y(y,z)=y,\qquad p_Z(y,z)=z.
\]
The diagram representing the projections and the component maps is
\[
\begin{tikzcd}[column sep=large, row sep=large]
 & X \arrow[dl, "h"'] \arrow[dr, "g"]
     \arrow[d, dashed, "{\beta(g,h)}"] & \\
 Z & Y\times Z \arrow[l, "p_Z"] \arrow[r, "p_Y"'] & Y
\end{tikzcd}
\]
For $f:X\to Y\times Z$ we define
\[
  \alpha(f)=(p_Y\circ f,\,p_Z\circ f).
\]
Conversely, for $g:X\to Y$ and $h:X\to Z$, set
\[
  \beta(g,h)(x)=(g(x),h(x)),\qquad x\in X.
\]

\medskip
\noindent\emph{Aside in the manuscript.} Fix $a\in Z$ and consider
\[
 \iota_Y:Y\longrightarrow Y\times Z,\qquad
 \iota_Y(y)=(y,a).
\]

\medskip
\noindent We prove first that $\alpha\circ\beta=\operatorname{id}$.
For every $g:X\to Y$ and $h:X\to Z$,
\[
  \alpha\bigl(\beta(g,h)\bigr)=(g,h).
\]
Indeed,
\[
\begin{aligned}
 \bigl(p_Y\circ\beta(g,h)\bigr)(x)
   &=p_Y\bigl(\beta(g,h)(x)\bigr) \\
   &=p_Y(g(x),h(x))=g(x),
 \qquad \forall x\in X,\\[3pt]
 \bigl(p_Z\circ\beta(g,h)\bigr)(x)
   &=p_Z\bigl(\beta(g,h)(x)\bigr) \\
   &=p_Z(g(x),h(x))=h(x),
 \qquad \forall x\in X.
\end{aligned}
\]
Consequently,
\[
  (\alpha\circ\beta)(g,h)
    =(p_Y\circ\beta(g,h),\,p_Z\circ\beta(g,h))
    =(g,h).
\]

\medskip
\noindent Conversely, we show that $\beta\circ\alpha=\operatorname{id}$.
Let $f:X\to Y\times Z$. Then
\[
 (\beta\circ\alpha)(f)=f.
\]
For $x\in X$, write $f(x)=(y,z)$. We have
\[
\begin{aligned}
 \beta\bigl(\alpha(f)\bigr)(x)
   &=\beta(p_Y\circ f,\,p_Z\circ f)(x)\\
   &=\bigl((p_Y\circ f)(x),\,(p_Z\circ f)(x)\bigr)\\
   &=\bigl(p_Y(y,z),\,p_Z(y,z)\bigr)\\
   &=(y,z)=f(x).
\end{aligned}
\]
Thus the two correspondences are mutually inverse.
; no preamble or forced page break.
% EDITORIAL NOTES:
% - This page continues Problem 2(b) from Seminar 3 - Page 1.
% - The handwritten notation varies between P_y/P_z and p_y/p_z;
%   the projection maps are typeset consistently as p_Y and p_Z.
% - The diagram below rearranges the handwritten projection sketch for
%   clarity; the two component maps and the projections are retained.
% - A short aside seems to read: fix a in Z and define i_Y(y)=(y,a).
%   This is included, although it is not used in the following proof.
% - Crossed-out calculations near the two composition identities are
%   omitted; all legible final calculations are retained.
% - The manuscript occasionally writes an identity without a domain;
%   the domains are made explicit in the displayed proof.

\subsection*{Seminar 3 - Page 2}

\noindent\textbf{(b) Continued.} We have the correspondence
\[
  \operatorname{Hom}(X,Y\times Z)
  \;\overset{\alpha}{\underset{\beta}{\rightleftarrows}}\;
  \operatorname{Hom}(X,Y)\times\operatorname{Hom}(X,Z).
\]
Let $p_Y:Y\times Z\to Y$ and $p_Z:Y\times Z\to Z$ be the
canonical projections:
\[
 p_Y(y,z)=y,\qquad p_Z(y,z)=z.
\]
The diagram representing the projections and the component maps is
\[
\begin{tikzcd}[column sep=large, row sep=large]
 & X \arrow[dl, "h"'] \arrow[dr, "g"]
     \arrow[d, dashed, "{\beta(g,h)}"] & \\
 Z & Y\times Z \arrow[l, "p_Z"] \arrow[r, "p_Y"'] & Y
\end{tikzcd}
\]
For $f:X\to Y\times Z$ we define
\[
  \alpha(f)=(p_Y\circ f,\,p_Z\circ f).
\]
Conversely, for $g:X\to Y$ and $h:X\to Z$, set
\[
  \beta(g,h)(x)=(g(x),h(x)),\qquad x\in X.
\]

\medskip
\noindent\emph{Aside in the manuscript.} Fix $a\in Z$ and consider
\[
 \iota_Y:Y\longrightarrow Y\times Z,\qquad
 \iota_Y(y)=(y,a).
\]

\medskip
\noindent We prove first that $\alpha\circ\beta=\operatorname{id}$.
For every $g:X\to Y$ and $h:X\to Z$,
\[
  \alpha\bigl(\beta(g,h)\bigr)=(g,h).
\]
Indeed,
\[
\begin{aligned}
 \bigl(p_Y\circ\beta(g,h)\bigr)(x)
   &=p_Y\bigl(\beta(g,h)(x)\bigr) \\
   &=p_Y(g(x),h(x))=g(x),
 \qquad \forall x\in X,\\[3pt]
 \bigl(p_Z\circ\beta(g,h)\bigr)(x)
   &=p_Z\bigl(\beta(g,h)(x)\bigr) \\
   &=p_Z(g(x),h(x))=h(x),
 \qquad \forall x\in X.
\end{aligned}
\]
Consequently,
\[
  (\alpha\circ\beta)(g,h)
    =(p_Y\circ\beta(g,h),\,p_Z\circ\beta(g,h))
    =(g,h).
\]

\medskip
\noindent Conversely, we show that $\beta\circ\alpha=\operatorname{id}$.
Let $f:X\to Y\times Z$. Then
\[
 (\beta\circ\alpha)(f)=f.
\]
For $x\in X$, write $f(x)=(y,z)$. We have
\[
\begin{aligned}
 \beta\bigl(\alpha(f)\bigr)(x)
   &=\beta(p_Y\circ f,\,p_Z\circ f)(x)\\
   &=\bigl((p_Y\circ f)(x),\,(p_Z\circ f)(x)\bigr)\\
   &=\bigl(p_Y(y,z),\,p_Z(y,z)\bigr)\\
   &=(y,z)=f(x).
\end{aligned}
\]
Thus the two correspondences are mutually inverse.
